\section{Adaptación entre la arquitectura del modelo y el aprendizaje por refuerzo}

\subsection{El algoritmo se adapta a la arquitectura vs. la arquitectura se adapta al algoritmo}

Cui Ganqu (崔干曲) propuso un cambio de perspectiva interesante: normalmente consideramos «hacer que el algoritmo se adapte a la arquitectura del modelo», pero pocas veces se piensa a la inversa, «hacer que la arquitectura se adapte al algoritmo». Sostiene que el diseño de la arquitectura está impulsado sobre todo por la infraestructura y la eficiencia del hardware (eficiencia de entrenamiento, eficiencia de inferencia, soporte multimodal, procesamiento de contexto largo), y que el algoritmo de RL ocupa una posición relativamente «posterior»: se adapta después de que la arquitectura ya está definida.

\begin{warningbox}{Los nuevos retos que trae la arquitectura híbrida}
Entrenar RL sobre una arquitectura híbrida (como un modelo mixto de Mamba y Attention) es mucho más complejo que hacerlo sobre una arquitectura de Full Attention. Las arquitecturas híbridas involucran operaciones de modificación y reversión del estado, y calcular estas operaciones de forma eficiente no es trivial. Esto también implica que existe una compensación entre el entrenamiento de RL y el diseño de la arquitectura: quienes investigan RL necesitan coordinarse con el equipo de arquitectura para asegurar que las arquitecturas nuevas sean favorables para el entrenamiento de RL posterior.
\end{warningbox}

\subsection{Ser favorable para RL es ser favorable para la inferencia}

Hu Jian (胡建), desde la perspectiva de la infraestructura, dio un juicio certero:

\begin{importantbox}{Arquitectura de modelo favorable para RL = arquitectura de modelo favorable para la inferencia}
En los escenarios actuales de contexto largo y salidas de pensamiento largo, la mayor parte del cuello de botella de cómputo del entrenamiento de RL se concentra en la parte de generación (Generation/Rollout). Por lo tanto, mejorar la eficiencia del entrenamiento de RL equivale, en esencia, a mejorar la eficiencia de la inferencia. Las innovaciones de arquitectura favorables para la inferencia incluyen:
\begin{itemize}
    \item \textbf{Linear Attention / Mamba}: reduce la complejidad computacional de las secuencias largas
    \item \textbf{Qianwen 3 Next}: innovaciones de arquitectura en cuanto a favorabilidad para la inferencia
    \item \textbf{MTP (Multi-Token Prediction)}: propuesto por DeepSeek, mejora de forma notable la eficiencia de la inferencia
\end{itemize}
\end{importantbox}

Hu Jian (胡建) también planteó una idea prospectiva: si en el futuro el \textbf{Diffusion LLM} sustituyera al paradigma autorregresivo, el cuello de botella de todo el entrenamiento de RL podría desplazarse de la optimización de la inferencia a la optimización del entrenamiento, lo cual sería un cambio radical.

\subsection{Dos perspectivas sobre la relación entre RL y los modelos grandes}

Zheng Chujie (郑楚杰) propuso dos perspectivas contrapuestas:

\textbf{Perspectiva uno: RL como auxiliar del modelo grande.} Primero se fija la arquitectura del modelo (MoE, MTP y otras técnicas de aceleración de la inferencia), y luego se adapta el algoritmo de RL sobre esa base. Esta es la práctica dominante actual en la industria.

\textbf{Perspectiva dos: el modelo grande como auxiliar de RL.} Si nuestro objetivo final es alcanzar la AGI mediante RL, entonces el modelo grande es solo un medio para aportar conocimiento previo. Desde esta perspectiva, debería usarse la arquitectura de modelo más favorable para el entrenamiento de RL (como un Dense Model, sin técnicas de aceleración como MTP). Pero en la práctica esta manera de proceder tiene pocos adeptos debido a su costo demasiado alto.

\begin{knowledgebox}{Los cambios en la arquitectura del modelo afectan las propiedades de RL}
Una observación importante de Zheng Chujie (郑楚杰): los cambios en la arquitectura del modelo (así como en el preentrenamiento del modelo base) modifican en cadena las propiedades del entrenamiento de RL posterior. Por eso, el problema de RL y el problema de la arquitectura no pueden verse por separado: cada generación de modelos requiere mejoras de RL específicas para sus propios problemas.
\end{knowledgebox}

\subsection{La dificultad del entrenamiento de RL en la arquitectura MoE}

Li Yingru (李英儒), desde la perspectiva de la teoría de la optimización, explicó por qué los modelos MoE son difíciles de entrenar con RL:

\begin{warningbox}{La raíz teórica de por qué MoE es difícil de entrenar}
\begin{itemize}
    \item La selectividad discreta de la arquitectura MoE (la selección de Expert por parte del Router) dificulta la optimización desde su base
    \item Desde la perspectiva de los algoritmos Off-Policy, MoE en realidad \textbf{reduce la Trust Region}: un cambio en la selección de Expert provoca un cambio drástico en el comportamiento del modelo
    \item Técnicas existentes como Routing Replay están precisamente resolviendo este problema de reducción de la Trust Region
    \item Este problema no es exclusivo de RL: MoE ya presentaba un problema similar de balance de enrutamiento durante el preentrenamiento (Pretraining)
\end{itemize}
\end{warningbox}

Li Yingru (李英儒) se mostró optimista respecto al futuro: podrían surgir nuevas arquitecturas Sparse que conserven la ventaja de eficiencia de la activación Sparse y, a la vez, tengan mejores propiedades de optimización (por ejemplo, aproximar mejor el gradiente en el Router).

\subsection{Resumen de la sección}

Existe una relación de acoplamiento profundo entre la arquitectura del modelo y el entrenamiento de RL. Actualmente, el paradigma dominante en la industria es «arquitectura primero, RL se adapta». Ser favorable para RL equivale, en esencia, a ser favorable para la inferencia: las innovaciones a nivel de arquitectura (Linear Attention, MTP, etcétera) benefician directamente la eficiencia del entrenamiento de RL. La arquitectura MoE trae consigo un problema particular de reducción de la Trust Region para RL, que requiere del diseño conjunto entre el algoritmo y la arquitectura para resolverse.
